%% Macros de courbes ajoutées pour la 1ère hybridation (données du manuscrit : colonnes h, errL2).
%% S'appuient sur l'environnement cvgh de convergences_h.tex.

% 1D : FEM P1 (plein), FEM P2 (plein), Add P1 (tireté), Mult M=3 et M=100 (pointillés) ; #1 dossier des CSV, #2 numéro du paramètre, #3 position des pentes
\newcommand{\cvgOneDAddMult}[3]{
    \begin{cvgh}[xticklabel style={rotate=20, anchor=east, yshift=-1pt}]{#1/FEM_case1_v1_param#2_degree1.csv}{3}{2,3}{#3}{$L^2$}{false}
        \addlegendentry{\,FEM $\mathbb{P}_1$\;}
        \addlegendentry{\,FEM $\mathbb{P}_2$\;}
        \addlegendentry{\,Add\;}
        \addlegendentry{\,Mult $M=3$\;}
        \addlegendentry{\,Mult $M=100$\;}

        \addplot [style={solid}, mark=square*, mark size=2, color=femcolor, line width=0.8pt]
        table [x=h, y=errL2, col sep=comma] {#1/FEM_case1_v1_param#2_degree1.csv};
        \addplot [style={solid}, mark=*, mark size=2, color=femcolor, line width=0.8pt]
        table [x=h, y=errL2, col sep=comma] {#1/FEM_case1_v1_param#2_degree2.csv};
        \addplot [style={dashed}, mark=square*, mark size=2, color=addcolor, line width=0.8pt]
        table [x=h, y=errL2, col sep=comma] {#1/Corr_case1_v1_param#2_degree1.csv};
        \addplot [style={densely dotted}, mark=square*, mark size=2, color=multcolor3, line width=1pt]
        table [x=h, y=errL2, col sep=comma] {#1/Mult_case1_v1_param#2_degree1_M3.0.csv};
        \addplot [style={densely dotted}, mark=square*, mark size=2, mark options={fill=white}, color=multcolor100, line width=1pt]
        table [x=h, y=errL2, col sep=comma] {#1/Mult_case1_v1_param#2_degree1_M100.0.csv};
    \end{cvgh}
}

% 1D, sans multiplicatif : FEM P1 et P2 (plein), Add P1 (tireté) ; mêmes arguments que \cvgOneDAddMult
\newcommand{\cvgOneDAdd}[3]{
    \begin{cvgh}[xticklabel style={rotate=20, anchor=east, yshift=-1pt}]{#1/FEM_case1_v1_param#2_degree1.csv}{3}{2,3}{#3}{$L^2$}{false}
        \addlegendentry{\,FEM $\mathbb{P}_1$\;}
        \addlegendentry{\,FEM $\mathbb{P}_2$\;}
        \addlegendentry{\,Add\;}

        \addplot [style={solid}, mark=square*, mark size=2, color=femcolor, line width=0.8pt]
        table [x=h, y=errL2, col sep=comma] {#1/FEM_case1_v1_param#2_degree1.csv};
        \addplot [style={solid}, mark=*, mark size=2, color=femcolor, line width=0.8pt]
        table [x=h, y=errL2, col sep=comma] {#1/FEM_case1_v1_param#2_degree2.csv};
        \addplot [style={dashed}, mark=square*, mark size=2, color=addcolor, line width=0.8pt]
        table [x=h, y=errL2, col sep=comma] {#1/Corr_case1_v1_param#2_degree1.csv};
    \end{cvgh}
}

% 1D (k=1, mu^(1)) : gain 1/C_gain de l'approche multiplicative en fonction de M (pointillé) contre le gain additif constant (tireté) ;
% #1 CSV (colonnes M, Cgainmult), #2 M maximal.
% Gain additif 1/0.02940315944770827 = 34.01 (script plot_gain_comparison.py du manuscrit, figures/efup/testcase1D/lap1D/comp_add_mult)
\newcommand{\cvgGainConstants}[2]{
    \begin{tikzpicture}
        \begin{semilogxaxis}[
                smaller labels,
                axis lines=left,
                enlarge x limits=false,
                xmin=0.5, xmax=#2, ymin=0, ymax=38,
                xlabel={$M$},
                xlabel style={at={(ticklabel* cs:1.01)},anchor=south west},
                ytick={0,10,20,30},
                width=\plotwidth, height=\plotheight,
                grid=major,
                xminorticks=false, yminorticks=false,
                legend style={font=\footnotesize, at={(0.97,0.05)}, anchor=south east, draw=none, fill=white, fill opacity=0.8, text opacity=1},
                legend image post style={mark options={solid, scale=1}},
            ]
            \addplot [dashed, color=addcolor, line width=1.2pt] coordinates {(0.5,34.01) (#2,34.01)};
            \addlegendentry{\,Add\;}
            \addplot [densely dotted, color=multcolor3, line width=1.4pt, each nth point=25, filter discard warning=false, unbounded coords=discard]
            table [x=M, y expr=1/\thisrow{Cgainmult}, col sep=space] {#1};
            \addlegendentry{\,Mult\;}
        \end{semilogxaxis}
    \end{tikzpicture}
}

% 3D (k=1) : FEM et Add, lus depuis les CSV du manuscrit (colonne errL2) ; #1 FEM, #2 Add, #3 position des pentes, #4 trait horizontal e (true/false)
\newcommand{\cvgFEMCorrOnedegErr}[4]{
    \begin{cvgh}[xticklabel style={rotate=20, anchor=east, yshift=-1pt}]{#1}{3}{2}{#3}{$L^2$}{#4}
        \addlegendentry{\,FEM $\mathbb{P}_1$\;}
        \addlegendentry{\,Add $\mathbb{P}_1$\;}
        \addplot [style={solid}, mark=square*, mark size=2, color=femcolor, line width=0.8pt]
        table [x=h, y=errL2, col sep=comma] {#1};
        \addplot [style={dashed}, mark=square*, mark size=2, color=addcolor, line width=0.8pt]
        table [x=h, y=errL2, col sep=comma] {#2};
    \end{cvgh}
}
